\subsection{Siêu phẳng phân tách}
\label{section:separate_hyper}


Hình \ref{fig:separate_hyper_example} minh họa 20 điểm dữ liệu trong hai lớp trong $\R^2$. Các dữ liệu này có thể được phân tách bởi một đường biên tuyến tính. Trong hình (đường màu xanh) là hai trong số vô số siêu phẳng phân tách có thể có. Đường màu cam là nghiệm bình phương tối thiểu của bài toán, thu được bằng cách hồi quy đáp ứng $-1/1$ của $Y$ trên $X$ (có hệ số chặn); đường này được cho bởi
\begin{equation}
\{x: \hat{\beta}_0 + \hat{\beta}_1 x_1 + \hat{\beta}_2 x_2 = 0\}
\label{eq:least_squares_boundary}
\end{equation}
Nghiệm bình phương tối thiểu này không thực hiện tốt việc phân tách các điểm và mắc một lỗi.